% =====================================================================
% Тезисы для конференции «Научный телеграф» (Центральный Университет /
% Т-Банк, май 2026). Стиль наследуется от paper/main.tex (REVTeX 4-2),
% но без twocolumn — тезисы на 1.5–2 страницы одной колонкой.
% =====================================================================
\documentclass[%
  aps,pre,
  superscriptaddress,
  longbibliography,nofootinbib,
  showkeys,
  amsmath,amssymb,floatfix
]{revtex4-2}

% --- Encoding & languages — Cyrillic-aware ---
\usepackage[T1,T2A]{fontenc}
\usepackage[utf8]{inputenc}
\usepackage[english,russian]{babel}
\usepackage{cmap}
\usepackage{lmodern}
\usepackage{microtype}

% --- Font setup: hybrid Times-math + CM-Cyrillic ---
\usepackage{mathptmx}
\renewcommand{\rmdefault}{cmr}
\DeclareMathSizes{10}{9}{6.5}{5}

% --- Math + symbols ---
\usepackage{amsmath,amssymb,amsfonts,mathtools}

% --- Floats / tables / figures ---
\usepackage{graphicx}
\usepackage{booktabs}
\setlength{\textfloatsep}{8pt plus 2pt minus 2pt}
\setlength{\floatsep}{8pt plus 2pt minus 2pt}
\setlength{\intextsep}{8pt plus 2pt minus 2pt}
\setlength{\abovecaptionskip}{4pt plus 1pt minus 1pt}
\setlength{\belowcaptionskip}{0pt}

% --- TikZ for the architecture diagram ---
\usepackage{tikz}
\usetikzlibrary{positioning,arrows.meta,calc,fit,shapes.geometric,
                automata,matrix,decorations.pathreplacing,patterns}
\tikzset{
  >={Latex[length=2mm,width=1.5mm]},
  block/.style ={rectangle, draw=black, thick, rounded corners=2pt,
                 minimum height=8mm, align=center, font=\small},
  arrow/.style ={-Stealth, thick},
}

% --- Hyperref + numeric square-bracket citations ---
\usepackage{hyperref}
\setcitestyle{numbers,square,sort&compress}

% --- Page footer with conference banner (same as main paper) ---
\usepackage{fancyhdr}
\newcommand{\confbanner}{%
  \itshape\footnotesize
  Мультидисциплинарная молодежная конференция
  Центрального университета \guillemotleft Научный телеграф\guillemotright}
\fancypagestyle{withbanner}{%
  \fancyhf{}%
  \fancyfoot[C]{\confbanner}%
  \fancyfoot[R]{\thepage}%
  \renewcommand{\headrulewidth}{0pt}%
  \renewcommand{\footrulewidth}{0pt}%
}
\pagestyle{withbanner}

% --- Math shortcuts (used in main paper) ---
\providecommand{\Real}{\mathbb{R}}
\providecommand{\Expect}{\mathbb{E}}
\providecommand{\Prb}{\mathbb{P}}
\DeclareMathOperator*{\argmaxop}{arg\,max}
\DeclareMathOperator*{\argminop}{arg\,min}

% --- Cross-reference helpers ---
\newcommand{\eqlab}[1]{\label{eq:#1}}
\newcommand{\seclab}[1]{\label{sec:#1}}
\newcommand{\figlab}[1]{\label{fig:#1}}
\newcommand{\Eqref}[1]{уравн.~\eqref{eq:#1}}
\newcommand{\Figref}[1]{РИС.~\ref{fig:#1}}

% --- Section heading macro: same look as main paper ---
\newcounter{musecct}
\renewcommand{\themusecct}{\Roman{musecct}}
\makeatletter
\newcommand{\sect}[2]{%
  \par\addvspace{0.35\baselineskip}%
  \refstepcounter{musecct}%
  \begingroup
    \centering
    \normalsize\bfseries
    \MakeUppercase{#1}%
    \par
  \endgroup
  \seclab{#2}%
  \nopagebreak
  \addvspace{0.12\baselineskip}%
  \@afterindentfalse\@afterheading
  \ignorespaces}
\makeatother

\renewcommand{\thetable}{\arabic{table}}
\addto\captionsrussian{%
  \renewcommand{\figurename}{РИС.}%
  \renewcommand{\tablename}{ТАБЛ.}%
}

% Russian "Ключевые слова:" label
\makeatletter
\renewcommand{\@keys@name}{{\small\bfseries Ключевые слова:}\space}
\makeatother

% Author-name helpers
\providecommand{\firstname}[1]{#1}
\providecommand{\surname}[1]{#1}
\makeatletter
\renewcommand{\andname}{}
\def\@listand{\@ifnum{\@tempcnta=\tw@}{,\space}{}}
\makeatother

% =====================================================================
\begin{document}
\selectlanguage{russian}

\keywords{отслеживание партитуры, скрытые марковские модели,
          on-line time warping, обучение с подкреплением,
          опережающее предсказание темпа,
          виртуальный оркестр}

\title{\MakeUppercase{Теоретические основы отслеживания партитуры в
       реальном времени: от скрытых марковских моделей к обучению с
       подкреплением}}

\author{\firstname{Н.~А.}~\surname{Новицкий}}
\email{n.novitskiy@edu.centraluniversity.ru}
\affiliation{\textit{АНО Центральный Университет, бывш. НИУ ВШЭ}}

\author{\firstname{Н.~С.}~\surname{Борисов}}
\email{n.borisov@edu.centraluniversity.ru}
\affiliation{\textit{АНО Центральный Университет}}

\begin{abstract}
\textbf{Аннотация.}\space
В работе проведено систематическое теоретическое изложение моделей
и алгоритмов отслеживания партитуры в реальном времени с акцентом
на современное направление~--- применение обучения с подкреплением.
Рассмотрены классические алгоритмы динамического программирования,
скрытые марковские модели и их полу-марковские расширения, а также
пять известных в литературе RL-применений к музыкальному
аккомпанементу. Сформулирована теоретическая постановка гибридного
подхода, в котором RL-агент действует как тонкий слой опережающего
предсказания темпа поверх классического вероятностного трекера, а
не заменяет его целиком. Численные эксперименты не приводятся;
работа имеет обзорный и теоретический характер.
\end{abstract}

\maketitle
\thispagestyle{withbanner}

\sect{Актуальность и цель}{intro}
%
Системы автоматического аккомпанемента нужны учащимся
консерваторий и концертирующим солистам как замена живого
оркестра при подготовке концертных произведений. Технически
задача сводится к отслеживанию партитуры в реальном времени
(\textit{score following})~\cite{Dannenberg1984}: по поступающему
каузально потоку наблюдений и заранее известной партитуре
требуется оценивать текущую позицию исполнителя с задержкой не
более $20$\,мс~--- порога восприятия рассинхронизации в
музыкальном ансамбле~\cite{Cont2010}. За четыре десятилетия в
области сменилось несколько парадигм: классическое динамическое
программирование~\cite{SakoeChiba1978,Dixon2005}, скрытые
марковские модели и их полу-марковские
расширения~\cite{Rabiner1989,Cont2010,Nakamura2015}, современные
нейросетевые подходы~\cite{Henkel2021} и относительно новое
направление, формализующее задачу как обучение с
подкреплением~\cite{Dorfer2018RL,Peter2023,Wu2024ReaLchords}.
Появление работ~\cite{Dorfer2018RL,Peter2023,Wu2024ReaLchords}
делает актуальным систематический теоретический разбор всего
семейства методов. Цель работы~--- дать математически выверенный,
прослеживаемый по литературе фундамент и явно сформулировать,
чего в существующих RL-подходах не сделано.

\sect{Основная идея}{idea}
%
Из обзора работ~\cite{Dorfer2018RL,Peter2023,Jiang2020RLDuet,
Wu2024ReaLchords,Jaques2017SeqTutor} следует, что чистая
RL-замена классического трекера~\cite{Dorfer2018RL} требует
тщательно размеченных данных уровня <<позиция-в-каждый-момент>> и
не использует априорную структуру партитуры (повторы, скачки,
явные длительности нот), естественно описываемую полу-марковскими
моделями. С другой стороны, классические трекеры~\cite{Cont2010}
обладают зрелой математической базой, но их
\textit{anticipation}-механизм зашит в виде явных байесовских
обновлений и неадаптивен к стилю солиста. Предлагается
\textit{гибридная} архитектура (\Figref{rl-agent}): классический
вероятностный трекер сохраняется как ядро системы, а RL-агент
действует как \textit{тонкий слой опережающего предсказания
темпа} на горизонте $200$--$500$\,мс с состоянием
$s_t = (\hat\alpha_t, \tau_{t-K:t}, e_{t-K:t}, \hat\varphi(t)/N)$
и непрерывным действием
$a_t \in [0{,}5;\,2{,}0]$~--- мультипликативным темповым
коэффициентом. Вознаграждение
\begin{equation}
r_t = -|t^{\text{render}}_t - t^{\text{perf}}_t|
      - \lambda\,\mathcal{L}_{\text{align}}(t)
      - \mu\,(a_t - a_{t-1})^2
\eqlab{reward-thesis}
\end{equation}
сочетает временное рассогласование, потерю выравнивания базового
трекера и штраф за резкие изменения темпа. Политика
$\pi_\theta(a\mid s)$ обучается двустадийно: сначала
behavioural cloning по выровненным учителям, затем PPO~\cite{Schulman2017PPO}.

\par\addvspace{2pt}
{\centering
\refstepcounter{figure}\figlab{rl-agent}%
\resizebox{0.55\linewidth}{!}{% =====================================================================
% TikZ figure: RL anticipation agent (BLACK & WHITE)
% Compact layout: 3 boxes top, 2 boxes bottom; fits into one column.
% Used in: sections/07-rl-anticipation.tex
% =====================================================================
\begin{tikzpicture}[
    node distance = 4mm and 3mm,
    every node/.append style={font=\scriptsize},
    block/.style ={rectangle, draw=black, thick, rounded corners=2pt,
                   minimum height=9mm, minimum width=20mm,
                   align=center, font=\scriptsize},
    arrow/.style ={-Stealth, thick}]

  \node[block]                  (track)  {Базовый\\трекер};
  \node[block, right=of track]  (state)  {Состояние\\$s_t$};
  \node[block, right=of state]  (policy) {Политика\\$\pi_\theta$};

  \node[block, below=7mm of policy] (env)    {Среда\\(рендерер)};
  \node[block, left=of env]         (reward) {Награда\\$r_t$};

  \draw[arrow] (track)  -- (state);
  \draw[arrow] (state)  -- (policy);
  \draw[arrow] (policy) -- node[right, font=\tiny] {$a_t$} (env);
  \draw[arrow] (env)    -- (reward);
  \draw[arrow] (reward.west) -| (track.south);

\end{tikzpicture}
}\par
\vspace{2pt}
\footnotesize
\figurename~\thefigure. Схема предлагаемой гибридной архитектуры:
базовый вероятностный трекер выдаёт состояние~$s_t$, RL-агент~---
темповый коэффициент~$a_t$ на горизонте $200$--$500$\,мс,
рендерер и симулятор солиста замыкают обратную связь через
вознаграждение $r_t$.\par}
\addvspace{4pt}

\sect{Положение и ценность}{position}
%
В рамках известных авторам работ предлагаемая схема не попадает
ни в одну из четырёх групп: (i)~RL-агент \textit{заменяет} трекер
целиком~\cite{Dorfer2018RL}; (ii)~RL-агент в myopic-режиме
решает задачу онлайн-выравнивания символьного
входа~\cite{Peter2023}; (iii)~RL применяется к \textit{генерации}
нового музыкального материала~\cite{Jiang2020RLDuet,Wu2024ReaLchords};
(iv)~RL применяется к общему KL-контролю предобученной
последовательностной модели~\cite{Jaques2017SeqTutor}. Ни одна
из перечисленных работ не рассматривает RL-агента как отдельный
тонкий слой опережающего предсказания темпа поверх сохраняемого
классического трекера. Ожидаемые свойства предложенной
рамки~--- вероятностная интерпретируемость через
forward-распределение, возможность адаптации под индивидуального
солиста и разделение задач между трекером и агентом для
упрощения диагностики. Полная версия работы~---
$19$-страничная статья~\cite{Novitsky2026Paper}~--- содержит
подробное изложение математического аппарата с числовыми
примерами и обзор 26 работ.

\bibliographystyle{apsrev4-2}
\begin{thebibliography}{99}
\bibitem{Dannenberg1984}
\textit{Dannenberg R. B.} An On-Line Algorithm for Real-Time
Accompaniment // Proc. of the International Computer Music
Conference (ICMC). --- 1984. --- P. 193--198.

\bibitem{SakoeChiba1978}
\textit{Sakoe H., Chiba S.} Dynamic Programming Algorithm
Optimization for Spoken Word Recognition // IEEE Trans. on
Acoustics, Speech, and Signal Processing. --- 1978. ---
Vol. ASSP-26, No. 1. --- P. 43--49.

\bibitem{Dixon2005}
\textit{Dixon S.} Live Tracking of Musical Performances Using
On-Line Time Warping // Proc. of the 8th Int. Conf. on Digital
Audio Effects (DAFx). --- 2005. --- P. 92--97.

\bibitem{Rabiner1989}
\textit{Rabiner L. R.} A Tutorial on Hidden Markov Models and
Selected Applications in Speech Recognition // Proc. of the
IEEE. --- 1989. --- Vol. 77, No. 2. --- P. 257--286.

\bibitem{Cont2010}
\textit{Cont A.} A Coupled Duration-Focused Architecture for
Real-Time Music-to-Score Alignment // IEEE Trans. on Pattern
Analysis and Machine Intelligence. --- 2010. --- Vol. 32,
No. 6. --- P. 974--987.

\bibitem{Nakamura2015}
\textit{Nakamura T., Nakamura E., Sagayama S.} Real-Time
Audio-to-Score Alignment of Music Performances Containing
Errors and Arbitrary Repeats and Skips // IEEE/ACM TASLP. ---
2015. --- Vol. 23, No. 11. --- P. 1911--1925.

\bibitem{Henkel2021}
\textit{Henkel F., Widmer G.} Real-Time Music Following in Score
Sheet Images via Multi-Resolution Prediction // Frontiers in
Computer Science. --- 2021. --- Vol. 3. --- Article 718340.

\bibitem{Dorfer2018RL}
\textit{Dorfer M., Henkel F., Widmer G.} Learning to Listen,
Read, and Follow: Score Following as a Reinforcement Learning
Game // Proc. of ISMIR 2018. --- P. 784--791.

\bibitem{Peter2023}
\textit{Peter S. D.} Online Symbolic Music Alignment with
Offline Reinforcement Learning // Proc. of ISMIR 2023. ---
arXiv:2401.00466.

\bibitem{Jiang2020RLDuet}
\textit{Jiang N., Jin S., Duan Z., Zhang C.} RL-Duet: Online
Music Accompaniment Generation Using Deep Reinforcement
Learning // Proc. of AAAI. --- 2020. --- Vol. 34. ---
P. 710--718.

\bibitem{Wu2024ReaLchords}
\textit{Wu Y. et al.} Adaptive Accompaniment with ReaLchords //
Proc. of the 41st International Conference on Machine Learning
(ICML). --- 2024. --- PMLR 235.

\bibitem{Schulman2017PPO}
\textit{Schulman J., Wolski F., Dhariwal P., Radford A.,
Klimov O.} Proximal Policy Optimization Algorithms //
arXiv preprint arXiv:1707.06347. --- 2017.

\bibitem{Jaques2017SeqTutor}
\textit{Jaques N. et al.} Sequence Tutor: Conservative
Fine-Tuning of Sequence Generation Models with KL-control //
Proc. of ICML 2017. --- PMLR 70. --- P. 1645--1654.

\bibitem{Novitsky2026Paper}
\textit{Новицкий Н. А., Борисов Н. С.} Теоретические основы
отслеживания партитуры в реальном времени: от скрытых
марковских моделей и on-line time warping к обучению с
подкреплением. --- Полная версия работы. --- 2026.
\end{thebibliography}

\end{document}
